\documentclass[11pt]{article}
\usepackage[margin=1in]{geometry}
\usepackage{amsmath,amssymb,amsthm,hyperref,mathrsfs}
\newtheorem{lemma}{Lemma}\newtheorem{theorem}[lemma]{Theorem}\newtheorem{corollary}[lemma]{Corollary}\newtheorem{proposition}[lemma]{Proposition}
\title{A digit-weight normalization of every physical gap-six matrix}
\author{Internal proof; no external-review or priority claim}\date{18 November 2026}
\begin{document}\maketitle
Let $s(n)=s_7(n)$ be the base-seven digit sum, with $s(0)=0$.
The coefficient variable $v$ is formal and unweighted.
For $f=\sum_{n\ge0}f_n(v)t^n\in\mathbb Z_7[v][[t]]$,
write $f\in\mathcal F_j$ if every coefficient of $f_n(v)$
has valuation at least $(s(n)-j)/6$.
This definition is also used for integral coefficients in a finite
extension of $\mathbb Q_7$, with $v_7(7)=1$.
Statements about infinite sums below are coefficientwise in $t$;
the one geometric expansion in the proof is coefficientwise
$7$-adically convergent.

\begin{lemma}[Digit filtration]\label{lem:filtration}
For integers $i,j$ one has
$\mathcal F_i\mathcal F_j\subseteq\mathcal F_{i+j}$,
$7\mathcal F_j\subseteq\mathcal F_{j-6}$, and
$\mathcal F_i\subseteq\mathcal F_j$ if $i\le j$.
Among integral series, $\mathcal F_1$ is closed under seventh
powers. It is closed under coefficientwise limits whenever the
limits exist.
\end{lemma}
\begin{proof}
The first assertion uses $s(a+b)\le s(a)+s(b)$.
The other filtration rules follow directly from the definition.
In a seventh power, a mixed multinomial term has coefficient
divisible by seven. If its indices, with multiplicities, are
$n_1,\ldots,n_7$, its valuation is at least
\[
 1+\sum_i\frac{s(n_i)-1}{6}
 \ \ge\ \frac{s(\sum_i n_i)-1}{6}.
\]
A pure term is $f_n(v)^7t^{7n}$; since the coefficients are integral,
$7v_7(f_n)\ge v_7(f_n)$ and $s(7n)=s(n)$ give the same bound.
This argument applies to polynomial coefficients in $v$ using
their Gauss valuation. Limits preserve each coefficient bound.
\end{proof}

For $N\ge1$ put $c=N+5$ and $A=c-\theta_z+z$.
Let $\Pi_N(X)=\prod_{\ell=0}^{N-1}(X+6\ell)$.
For $Q\in\mathbb Z_7[X]$ define the normalized boundary components
by the finite identity
\[
 \frac{\Pi_N(A)Q(A)1}{6^NN!}
    =\sum_j b_{N,Q,j}z^j,\qquad
 V_{N,a}(Q;v)=a!\sum_{\ell\ge0}
            b_{N,Q,a+6\ell}7^\ell v^\ell,\quad 0\le a<6.
 \tag{1}
\]
The finite-difference lattice in
\href{GAP_SIX_BOUNDARY_COUPLED_FROBENIUS.tex}
{Lemma \texttt{lem:mahler}}
proves integrality. A fresh proof of the stronger bound follows.

\begin{theorem}[Uniform all-index row divisibility]\label{thm:content}
For every $N\ge1$, $0\le a<6$, and $Q\in\mathbb Z_7[X]$,
\[
 V_{N,a}(Q;v)\in7^{e_a(N)}\mathbb Z_7[v],
 \qquad e_a(N)=
       \left\lceil\frac{s(N)-a-1}{6}\right\rceil\ge0.
 \tag{2}
\]
The same conclusion holds for polynomials $Q$ with integral
coefficients in any extension of $\mathbb Q_7$, with the
coefficientwise valuation interpretation of (2).
No degree restriction on $Q$ is required.
\end{theorem}
\begin{proof}
Use the exact diagonal identity in
\href{GAP_SIX_TWO_POWER_PLUS_ONE_SFC.tex}
{Lemma \texttt{lem:diagonal}},
which is an identity for all diagonal indices:
\[
 \sum_{N\ge0}\frac{B_{N,N,6}(z)}{6^NN!}t^N
       =\kappa(t)\exp(zH(t)),\quad
 H=1-R,\quad R-R^7=t,\quad
 \kappa=\frac{6R}{7R^6-1}.
 \tag{3}
\]
The equation for $H$ is
\[
 H=\frac t6+\frac72H^2-\frac{35}6H^3
          +\frac{35}6H^4-\frac72H^5+\frac76H^6-\frac16H^7 .
 \tag{4}
\]
Its right side contracts in the $t$-adic topology on
$t\mathbb Z_7[[t]]$. Lemma~\ref{lem:filtration} shows that it
preserves $\mathcal F_1$: for $2\le j\le6$ the coefficient of
$H^j$ is seven-divisible and $j-6\le1$, while $H^7$ lies in
$\mathcal F_1$. Consequently $H\in\mathcal F_1$.

Use the integral Artin--Hasse series
$E_7(X)=\exp(\sum_{i\ge0}X^{7^i}/7^i)$, whose integrality is
proved in
\href{ARITHMETIC_FRONTIER_TERMINAL_ORDERS.tex}
{Lemma \texttt{lem:artin-hasse-closure}}.
Put $E_i=(7^i-1)/6$ and let $\mathscr A\in
t\mathbb Z_7[v][[t]]$ solve
\[
 \mathscr A+\sum_{i\ge1}7^{E_i-i}v^{E_i}
                   \mathscr A^{7^i}=H.                 \tag{5}
\]
All exponents $E_i-i$ are nonnegative. Iteration of (5) is
$t$-adically contracting. Iterated seventh-power closure proves
$\mathscr A\in\mathcal F_1$.

At the formal boundary $z^6=7v$, equation (5) gives the exact
formal identity $\exp(zH)=E_7(z\mathscr A)$.
Each fixed $t$-coefficient is a finite polynomial in $z$, so no
$p$-adic exponential is evaluated at a boundary point.
Write $E_7(X)=\sum_{j\ge0}e_jX^j$, where $e_j\in\mathbb Z_7$.
The $a$th normalized cyclic component is
\[
 a!\sum_{\ell\ge0}e_{a+6\ell}7^\ell v^\ell
                          \mathscr A^{a+6\ell}
                  \in\mathcal F_a.                    \tag{6}
\]
Every summand is in $\mathcal F_a$, by the product and seven
rules; the sum is $t$-adically well-defined.

Also
\[
 \kappa=-6\sum_{j\ge0}7^j(1-H)^{6j+1}\in\mathcal F_1.
 \tag{7}
\]
Indeed every binomial term $7^jH^k$ has
$k-6j\le1$. This expansion converges coefficientwise
$7$-adically. Equations (3) and (6)--(7) prove (2) for $Q=1$.

The exact raising identities, before reduction, are
\[
 \begin{split}
 V_{N,a}(XQ)&=(c-a-6vD_v)V_{N,a}(Q)+aV_{N,a-1}(Q),
                                  \quad 1\le a<6,\\
 V_{N,0}(XQ)&=(c-6vD_v)V_{N,0}(Q)
                         +\frac{7v}{5!}V_{N,5}(Q).
 \end{split}                                             \tag{8}
\]
They follow by applying $A$ to the finite polynomial in (1).
Now $e_{a-1}\ge e_a$ and
\[
 1+e_5=1+\left\lceil\frac{s(N)-6}{6}\right\rceil
       =\left\lceil\frac{s(N)}6\right\rceil\ge e_0 .
\]
Thus (8) preserves the six bounds, proving them for every
monomial $Q=X^j$ and then every integral $Q$ by linearity.
The extension-field assertion follows from the same finite
linear combination, without integer-valued binomial assertions
outside $\mathbb Z_7$.
\end{proof}

\begin{corollary}[A single cut in the normalized raising cycle]
\label{cor:cut}
Set $\widetilde V_{N,a}=7^{-e_a(N)}V_{N,a}$, and put
$a_*=(s(N)-1)\bmod6=(N-1)\bmod6$.
The reduction of (8) is an open six-row chain: all cyclic
off-diagonal links survive except the one whose target is $a_*$.
Explicitly, for $1\le a<6$,
\[
 \overline{\widetilde V_a(XQ)}
   =(\bar c-a+vD_v)\overline{\widetilde V_a(Q)}
       +a\,\epsilon_a\overline{\widetilde V_{a-1}(Q)},
 \quad
 \epsilon_a=\begin{cases}0&a=a_*,\\1&a\ne a_*,\end{cases}
 \tag{9}
\]
and
\[
 \overline{\widetilde V_0(XQ)}
 =(\bar c+vD_v)\overline{\widetilde V_0(Q)}
       +v\epsilon_0\overline{\widetilde V_5(Q)}.
 \tag{10}
\]
Thus the raising rule has only $42$ possibilities, indexed by
$N\bmod7$ and $N\bmod6$. Its actual six seed polynomials
$\overline{\widetilde V_a(1)}$ are not asserted to belong to a
finite set or to have full rank.
\end{corollary}
\begin{proof}
The exponents of seven multiplying the five internal links are
$e_{a-1}-e_a\in\{0,1\}$; the wrap link has exponent
$1+e_5-e_0\in\{0,1\}$. Their sum is one, and their unique
nonzero term has target $(s(N)-1)\bmod6$.
Finally $-6=1$ and $5!=1$ in $\mathbb F_7$.
\end{proof}

\begin{lemma}[A finite polynomial gauge and seven coupled eigensections]
\label{lem:gauge}
Let $\kappa=\bar N+5\in\mathbb F_7$, $t=(N-1)\bmod6$, and
work on polynomial or Laurent-polynomial columns over
$\overline{\mathbb F}_7$. Put $\delta=vD_v$,
\[
 \mathcal D=\operatorname{diag}(\delta+\kappa-a)_{a=0}^5.
\]
Let $\mathcal J$ have entries $\mathcal J_{a,a-1}=a$
for $1\le a<6$, and $\mathcal J_{0,5}=v$, except that the
entry with target row $t$ is omitted.
The raising operator in (9)--(10) is
$\mathcal L=\mathcal D+\mathcal J$. One has
\[
 \mathcal J^6=0,\qquad
 E=\sum_{j=0}^5\frac{\mathcal J^j}{j!}\in
              \operatorname{SL}_6(\mathbb F_7[v]),\qquad
 \mathcal L=E\mathcal D E^{-1},\qquad
 \mathcal L^7=\mathcal L.                              \tag{11}
\]
For the actual seed column
$s_N=(\overline{\widetilde V_{N,a}(1)})_{a=0}^5$ put
$t_N=E^{-1}s_N$. Define $t_{N,\lambda}$ by retaining in
row $a$ precisely the monomials $v^k$ of $t_N$ with
$k+\kappa-a=\lambda$ in $\mathbb F_7$.
Then the entire normalized five-column matrix has the exact
polynomial factorization
\[
 [s_N,\mathcal Ls_N,\ldots,\mathcal L^4s_N]
  =E\,[t_{N,\lambda}]_{\lambda\in\mathbb F_7}
                  [\lambda^j]_{\lambda\in\mathbb F_7,\ 0\le j\le4}.
                                                               \tag{12}
\]
Thus evaluation at a residue point preserves the rank if the
left factor $E$ is removed. All seven eigensections in the
remaining six-by-seven matrix are still coupled.
\end{lemma}
\begin{proof}
The single removed link makes $\mathcal J$ an open chain,
so $\mathcal J^6=0$; a simultaneous ordering along that chain
makes it strictly triangular. Consequently $E$ has determinant
one and inverse $\sum_{j=0}^5(-\mathcal J)^j/j!$.
For an internal link,
$[\mathcal D,\mathcal J]_{a,a-1}=-\mathcal J_{a,a-1}$.
For the wrap link the Euler derivative of $v$ and the diagonal
difference contribute $1+5=6=-1$. Thus
$[\mathcal D,\mathcal J]=-\mathcal J$, and the finite
conjugation identity gives
$E\mathcal D E^{-1}=\mathcal D+\mathcal J$.
Since $\delta^7=\delta$ on Laurent monomials and every
$\kappa-a$ belongs to $\mathbb F_7$, one has
$\mathcal D^7=\mathcal D$, hence (11).

The eigensection $t_{N,\lambda}$ is exactly the component
with $\mathcal D t_{N,\lambda}=\lambda t_{N,\lambda}$.
Summing these components gives $t_N$, so conjugation gives
(12). No derivative is applied on an evaluated quotient.
Distinct eigensections can cancel after evaluation; separate
vanishing of them is not asserted.
\end{proof}

\begin{proposition}[The normalized residue rank can still drop]
\label{prop:rank-obstruction}
For $N=8$ and $\bar v=6$, the normalized five-column matrix
has rank exactly four. Hence full rank of every normalized
residue matrix is not an available all-index conclusion.
\end{proposition}
\begin{proof}
Here $s(8)=2$ and $(e_0,\ldots,e_5)=(1,0,0,0,0,0)$.
Direct finite recurrence (1), or the defining auxiliary recurrence,
gives the exact reduced matrix
\[
 \begin{pmatrix}
 6&1&6&1&4\\
 4&5&5&3&2\\
 4&5&3&1&1\\
 0&5&1&4&4\\
 0&0&6&5&3\\
 0&0&0&2&0
 \end{pmatrix}.
\]
All six maximal minors vanish in $\mathbb F_7$, while the
upper left four-by-four minor is one.
This point does not produce a multiple physical root:
$Y=7t=-1/v$ reduces to one, and the actual physical polynomial
$H_8(Y)=7T_{8,8,6}(Y/7)$ satisfies
$H_8(1)\equiv21\pmod{49}$.
Its leading coefficient is a unit by
\href{PHYSICAL_PRIME_TOTAL_DEGREE_REDUCTION.tex}
{Lemma \texttt{lem:coefficients}}, so its entire residue cluster
at one is Eisenstein and therefore separable, by
\href{GAP_SIX_SMALL_OFFSET_FAMILIES.tex}
{Lemma \texttt{lem:local-cluster}}.
\end{proof}

For example, $s(N)=8$ forces a second division by seven in row
zero, while the other five rows are guaranteed one division.
The normalization and finite gauge apply to every physical
index. They do not prove that the actual eigensection matrix
in (12) has rank five on the required points, that its actual
seeds have finite Cartier closure, or that physical gap-six
polynomials are squarefree. The remaining problem is to exploit
their actual diagonal-series origin rather than replace them
by arbitrary integral seeds.
\end{document}
